\documentclass[dvipsnames, svgnames, x11names]{beamer}
\usepackage[UTF8, noindent]{ctexcap}  % 中文支持、无段前缩进

\usepackage{multicol}  % 双栏目录
\usepackage{setspace}  % 调整行距
\usepackage{listings}

% \usepackage{ctex}
\usepackage{tikz}
\usepackage{xcolor}

% 代码抄录的一系列设置
\lstset{
    basicstyle=\footnotesize\ttfamily, 
    flexiblecolumns, 
    numbers=left, 
    numbersep=5pt, 
    numberstyle=\tiny, 
    backgroundcolor=\color{white}, 
    frame=single, 
    breaklines=true, 
    showtabs=false, 
    showspaces=false, 
    showstringspaces=false, 
    keywordstyle=\color{teal}, 
    commentstyle=\color{blue}, 
    stringstyle=\color{red}, 
    numberstyle=\color{gray}, 
    tabsize=8, 
    breakatwhitespace=false, 
    postbreak=\mbox{\textcolor{violet}{$\hookrightarrow$}\space}
} % 代码抄录的一系列设置


\mode<presentation> {
	%set the slide theme
    \useoutertheme[footline=authorinstitutetitle,subsection=false]{smoothbars}
    \useinnertheme{circles}
  \setbeamertemplate{footline} 
  {%
    \begin{beamercolorbox}[colsep=1.5pt]{upper separation line foot}
    \end{beamercolorbox}
    \begin{beamercolorbox}[ht=2.5ex,dp=1.125ex,%
      leftskip=.3cm,rightskip=.3cm plus1fil]{author in head/foot}%
      \leavevmode{\usebeamerfont{author in head/foot}\insertshortauthor}%
      \hfill%
      {\usebeamerfont{institute in head/foot}\usebeamercolor[fg]{institute in head/foot}\insertshortinstitute}%
    \end{beamercolorbox}%
    \begin{beamercolorbox}[ht=2.5ex,dp=1.125ex,%
      leftskip=.3cm,rightskip=.3cm plus1fil]{title in head/foot}%
      {\usebeamerfont{title in head/foot}\insertshorttitle}%
      \hfill%
      {\usebeamerfont{frame number}\usebeamercolor[fg]{frame number}\insertframenumber~/~\inserttotalframenumber}
    \end{beamercolorbox}%
    \begin{beamercolorbox}[colsep=1.5pt]{lower separation line foot}
    \end{beamercolorbox}
  }
	% \setbeamercovered{transparent}	
	\hypersetup{pdfpagemode=FullScreen}

    
% \xdefinecolor{csust_bg}{rgb}{0.03125,0.40625,0.70703125}  %RGB 8,104,181
% \xdefinecolor{csust_bg}{rgb}{0.046875, 0.390625, 0.625}  %RGB 12,100,160
\xdefinecolor{csust_bg}{rgb}{0.10546875, 0.19921875, 0.54296875}  %RGB 27,51,139
\setbeamercolor{footline}{bg=csust_bg}
\setbeamercolor{frametitle}{bg=csust_bg,fg=white}
\setbeamercolor{title}{bg=csust_bg}
\setbeamerfont{frametitle}{size=\large}
%\setbeamertemplate{navigation symbols}{}
% \setbeamertemplate{bibliography item}[text]
\setbeamertemplate{caption}[numbered]

\setbeamercolor{palette primary}{use=structure,fg=white,bg=structure.fg}
\setbeamercolor{palette secondary}{use=structure,fg=white,bg=structure.fg!75!black}
\setbeamercolor{palette tertiary}{use=structure,fg=white,bg=structure.fg!50!black}
\setbeamercolor{palette quaternary}{fg=white,bg=structure.fg!50!black}
%\setbeamercolor*{sidebar}{use=structure,bg=structure.fg}
\setbeamercolor{titlelike}{parent=palette primary}

%% try
\setbeamercolor{block title}{bg=csust_bg,fg=white}
\setbeamercolor*{block title example}{use={normal text,example text},bg=white,fg=csust_bg}
\setbeamercolor{fine separation line}{}
\setbeamercolor{item projected}{fg=white}
\setbeamercolor{palette sidebar primary}{use=normal text,fg=normal text.fg}
\setbeamercolor{palette sidebar quaternary}{use=structure,fg=structure.fg}
\setbeamercolor{palette sidebar secondary}{use=structure,fg=structure.fg}
\setbeamercolor{palette sidebar tertiary}{use=normal text,fg=normal text.fg}
%\setbeamercolor{palette sidebar quaternary}{fg=white}
\setbeamercolor{section in sidebar}{fg=brown}
\setbeamercolor{section in sidebar shaded}{fg=grey}
\setbeamercolor{separation line}{}
\setbeamercolor{sidebar}{bg=csust_bg}
\setbeamercolor{sidebar}{parent=palette primary}
\setbeamercolor{structure}{fg=csust_bg}
\setbeamercolor{subsection in sidebar}{fg=brown}
\setbeamercolor{subsection in sidebar shaded}{fg=grey}
    \AtBeginSection[]{
        \begin{frame}[plain]
            \tableofcontents[sectionstyle=show/shaded,subsectionstyle=show/shaded/hide,subsubsectionstyle=show/shaded/hide]
        \end{frame}
    }
    \AtBeginSubsection[]{
        \begin{frame}[plain]
            \tableofcontents[sectionstyle=show/shaded,subsectionstyle=show/shaded/hide,subsubsectionstyle=show/shaded/hide]
        \end{frame}
    }
} % end mode <presentation>

\title{长沙理工大学毕业答辩幻灯片 Beamer 模板}
\subtitle{CSUST Graduation Defense Slides Beamer Template}
\author[你的名字]{\makebox[4em][s]{演讲人}：\makebox[4em][l]{你的名字}\\ 指导教师：某某老师}
\institute[长沙理工大学计算机与通信工程学院]{长沙理工大学计算机与通信工程学院\\ 软件工程18-x班}
\date{\the\year 年6月7日}
\titlegraphic{\includegraphics[width=3cm]{../figure/csustlogo_626by572.jpg}}

\subject{编译原理}
\keywords{编译器、C++语言}

\begin{document}

\begin{frame}[plain]
    \titlepage
\end{frame}

\begin{frame}[plain]
    \tableofcontents[sectionstyle=show,subsectionstyle=show/shaded/hide,subsubsectionstyle=show/shaded/hide]
\end{frame}

\section{第一节}
\subsection{第1.1小节}
\begin{frame}{大标题}{小标题}
    \onslide<+-> 前面加上 onslide+-。
    \onslide<+-> 就可以一帧帧显示：
    \begin{itemize}[<+->]
        \item 列表环境也可以。
        \item 一帧帧显示。
    \end{itemize}
\end{frame}

\begin{frame}
    插入图片
    \begin{center}
        \includegraphics[scale=.5]{../figure/csust_logo_and_name.pdf}
    \end{center}
    \begin{center}
        \includegraphics[scale=.5]{../figure/csust_logo_and_name2.pdf}
    \end{center}
\end{frame}

\section{第二节}
\subsection{左文右图}
\begin{frame}
    \begin{minipage}{.5\linewidth}
        左文右图。
    \end{minipage}
    \quad\quad
    \onslide<1->
    \begin{minipage}{.4\linewidth}
    \includegraphics[width=\linewidth]{../figure/csustlogo_2640by2170.jpg}
    \end{minipage}
\end{frame}


\subsection{左图右文}
\begin{frame}
    \begin{minipage}{.4\linewidth}
    \includegraphics[width=\linewidth]{../figure/csustlogo_626by572.jpg}
    \end{minipage}
    \quad
    \begin{minipage}{.55\linewidth}
        左图右文。
    \end{minipage}
\end{frame}

\subsection{数学环境}
\begin{frame}
    勾股定理公式：
    \[
        a^2+b^2=c^2.
    \]
    勾三股四弦五：
    \[
        3^2+4^2=5^2.
    \]
\end{frame}

\section{结语}
\subsection{注意事项}
\begin{frame}{注意事项}
    本模板需要经过\alert{两次 XeLaTeX 编译}才能正确生成目录！
\end{frame}

\subsection{本模板的不足}
\begin{frame}{本模板的不足}
    \begin{itemize}[<+->]
        \item 缺乏示例。
        \item 缺乏文档。
        \item 更多不足有待发现……
    \end{itemize}

    \onslide<+-> 因此 长沙理工大学毕业答辩 Beamer 模板需要一届又一届长理学子接力共同维护。
\end{frame}

\begin{frame}[plain]
    \begin{center}
        {\Huge Thanks!}
    \end{center}
\end{frame}

\end{document}